\begin{proposition}\label{push:src:D.4.7}
The function $D_{\rm part}$ has no critical point on the punctured
root neighbourhood $0<|\tau-i|\leq3/20000$, no critical point on the
coarse neighbourhood represented by $|\tau-\rho|\leq1/50$, and no
critical point where $\Im\tau\geq2$. At $\rho$ the two factors
$Z$ and $F=3ZP-2g_2$ of \ref{push:src:D.4.4} vanish, but the coarse gradient is
nonzero, as in \ref{push:src:D.4.5}. The earlier interval proof and its 400 local rectangles remain in
the S6\_09 proof compendium under this label. We give the consequence
of the full critical-point classification here.
\end{proposition}
\Uses{\cref{ari:C.8.15}; \cref{push:src:D.4.4,push:src:D.4.5}; \cref{ari:C.1.1}}
\begin{proof}

The two norms agree by their exact sigma--theta formula in
\ref{ari:C.1.1}. Hence, with $L=\log\mathcal N$,
\[
 D_{\rm part}=\frac6\pi\bigl(L-\log(2\sqrt3)\bigr).
\]
Theorem~\ref{ari:C.8.15} proves that the coarse differential
is nonzero away from $t=1$, including at $t=0$, where
$L_t(0)=7/384$. It gives
$(D_{\rm part})_t(0)=7/(64\pi)$. Its two exact root-critical
orbits are $i$ and $\rho$, with the factor vanishing specified
in \ref{push:src:D.4.4,push:src:D.4.5}.

The root-to-coarse map is unramified off those elliptic orbits.
The stated small disc about $i$ contains no other elliptic
representative, and removing its centre therefore removes its
only critical point. The coarse disc about $\rho$ has the
nonzero differential just calculated. Finally
\[
 \Im\frac{ai+b}{ci+d}=\frac1{c^2+d^2}\le1,\qquad
 \Im\frac{a\rho+b}{c\rho+d}
 =\frac{\Im\rho}{c^2+cd+d^2}\le1.
\]
Thus no elliptic-orbit critical point lies at height at least two.
These facts prove every claimed domain of nonvanishing.

\end{proof}
